\documentclass[11pt]{article}
\usepackage[margin=1in]{geometry}
\usepackage{amsmath,amssymb}
\usepackage[hidelinks]{hyperref}

\title{Literature Status: The $\ell_2\oplus\mathcal T^{(2)}$ UTAP Subquestion}
\author{Run \texttt{fa\_banach\_001}, lane 4}
\date{2026-06-29}

\begin{document}
\maketitle

\section*{Status}

\textbf{literature\_already\_answered, scoped subquestion only.}

The source paper is F. Albiac, J. L. Ansorena, and P. Wojtaszczyk,
\emph{Uniqueness of unconditional basis of $H_p(\mathbb T)\oplus\ell_2$ and
$H_p(\mathbb T)\oplus\mathcal T^{(2)}$ for $0<p<1$}, arXiv:2010.01501.
In the final discussion after Theorem 4.7, it lists as one of the main
remaining Bourgain--Casazza--Lindenstrauss--Tzafriri style cases whether
\[
  \ell_2\oplus \mathcal T^{(2)}
\]
has a UTAP unconditional basis.

The supporting paper is F. Albiac and J. L. Ansorena,
\emph{Uniqueness of unconditional basis of $\ell_2\oplus\mathcal T^{(2)}$},
arXiv:2012.06783.  This is a separate later paper and it explicitly knows the
provenance: its introduction says it is a direct continuation of the 2020
work, explains that $\ell_2\oplus\mathcal T^{(2)}$ lies outside the previous
methods, and then says that the paper fills this gap by proving, in particular,
that $\ell_2\oplus\mathcal T^{(2)}$ has a UTAP unconditional basis.  The
abstract states the same conclusion.

\section*{Why This Answers The Source Subquestion}

The source question is an exact yes/no UTAP question for the direct sum
$\ell_2\oplus\mathcal T^{(2)}$.  The supporting paper's title, abstract, and
main application all identify the same direct sum and prove the affirmative
UTAP conclusion.  Therefore that subquestion is already answered in the
literature.

\section*{Scope Limitations}

This packet does not claim a solution of the broad question whether
$X\oplus Y$ has a UTAP unconditional basis whenever both Banach spaces have
one.  It also does not settle the later Gowers-square question from
arXiv:2603.07324 asking whether $G^m$ has a unique unconditional basis for
$m\geq2$.  The existing local attempt
\texttt{attempts/2010.01501\_gowers\_square\_vector\_peak\_reduction.md}
remains the relevant no-promotion note for that route.

\section*{Evidence Checked}

Local cheap indexes showed no prior solution packet for arXiv:2010.01501, but
did show the Gowers-square failed attempt.  The parsed local sources for
arXiv:2010.01501, arXiv:2012.06783, and the later Gowers route source
arXiv:2603.07324 were inspected.

\end{document}
